\section{Weekly Summary}

\begin{tcolorbox}[colback=yellow!10,colframe=black!50,title=AI-Generated Content]
\noindent The summary below was written by Claude (Anthropic) from that week's regression outputs and series descriptions. It has not been reviewed or fact-checked by a human, and should be read as an automated, informal recap -- not analysis.
\end{tcolorbox}

Welcome back to the weekly edition of Two Random Data Series Walk Into A Bar, where we bravely regress things that have no business being in the same spreadsheet and then act shocked at the results. This week our tireless robot pulled seven pairs out of FRED's endless bin of numbers, and honestly, the lineup did not disappoint. We opened Sunday by asking whether Lithuania's "Discussion About Pandemics Index" -- a series so obscure its popularity score is literally the number one, as in one lonely soul -- has anything to do with the U.S. Treasury borrowing from the public. And, uh, kind of? The raw regression was significant, and remarkably it held up detrended too, which is the rare and unsettling case where I'm forced to stop dunking and just stare into the middle distance. Do Lithuanians chatter about pandemics whenever Uncle Sam floats more debt? No. Absolutely not. But the p-value doesn't know that, and that's the fun part.

Monday gave us coal-fired carbon dioxide emissions versus the price index for treating congenital anomalies, which is a sentence I never planned to write. Raw R-squared of about 0.61, and it actually survived detrending with an even punchier slope. Two things drifting together for decades usually means nothing, yet here the relationship refused to evaporate, so I'll grudgingly file it under "huh." Tuesday, meanwhile, was the classic trend mirage we live for: the market value of gross federal debt versus business equipment leases owned by finance companies. The raw regression insisted there was a link with a p-value with eleven zeros in it, but of course both series just spent decades marching upward like everything else denominated in dollars. Detrended, the relationship flipped sign and got even more "significant," which mostly tells you the whole thing is numerical noise wearing a nice suit.

Wednesday's pairing -- construction quits versus the work-from-home rate for women -- was the week's comedy relief. The raw regression was a spectacular flatline, an R-squared of 0.0004 and a p-value of 0.86, which is statistics for "these two have never met." Then, plot twist, detrending yanked out a genuinely significant negative relationship, proving that once you strip the trends away sometimes a ghost appears where there was nothing before. Thursday brought national wages and salaries against manufacturing employment in Ogden, Utah, a raw R-squared of 0.80 that shrank to a wobblier but still-present 0.40 detrended, so partial credit. And Friday delivered Midwest shelter costs versus the price of a whole chicken in the West, which posted a jaw-dropping raw R-squared of 0.94 -- both rent and poultry go up over time, film at eleven -- but still held on detrended, suggesting that when the cost of living climbs, so does the cost of dinner. Groundbreaking.

We closed Saturday on the week's most legitimately boring result, and I mean that as a compliment: two employment cost indices, one for private production workers and one for state and local government, correlating at a raw R-squared of 0.99 and holding strong at 0.83 detrended. This is the one pairing all week where the relationship is actually plausible, because compensation costs really do move together across sectors -- so naturally it's the least interesting to write about. As always, the mandatory disclaimer: these are randomly matched series, no theory, no causation, no rigor, just a machine flinging numbers at each other for our amusement. Spurious correlation is the house special here, not the exception, so please do not adjust your portfolio, your carbon policy, or your chicken budget based on anything above. See you next week for more statistically significant nonsense.
